\chapter{Datos}

\subsection{Fuentes de datos}

El pipeline se entrena y evalua con tres tipos de datos:

\begin{itemize}[leftmargin=*,itemsep=2pt]
  \item \textbf{Sinteticos generados por el gemelo}: constituyen el grueso
        de la evaluacion. Permiten controlar el catalogo de fallas y
        variar condiciones operativas (perfil de carga, clima, mix de
        clientes) sin riesgo.
  \item \textbf{Publicos del CEN}: medidas de frecuencia y tension del
        SEN, usadas como sanity check de la dinamica de fondo del gemelo.
  \item \textbf{Reales (futuro)}: a traves de alianza con Chilquinta o
        Engie, ya iniciadas conversaciones preliminares. No se usan en
        este trabajo por restricciones de confidencialidad.
\end{itemize}

\subsection{Gemelo digital: parametros}

La subestacion gemela es de distribucion 110/23~kV con dos
transformadores de 30~MVA en paralelo y cuatro alimentadores de
23~kV. Los parametros del modelo se resumen en la
tabla~\ref{tab:parametros}.

\begin{table}[ht]
\centering
\small
\caption{Parametros de la subestacion gemela.}
\label{tab:parametros}
\begin{tabular}{lcc}
\toprule
\textbf{Parametro} & \textbf{Valor} & \textbf{Unidad} \\
\midrule
$S_{nom}$ por trafo & 30 & MVA \\
$V_{AT}$ / $V_{BT}$ & 110 / 23 & kV \\
$X_{cc}$ por trafo & 0.12 & pu \\
$P_{Cu,nom}$ & 180 & kW \\
$P_{Fe}$ & 35 & kW \\
$\tau_{termico}$ & 90 & min \\
$V_{min}$ / $V_{max}$ & 0.93 / 1.07 & pu \\
$f_{min}$ / $f_{max}$ & 49.5 / 50.5 & Hz \\
$P_{nom}$ por linea (L1..L4) & 6.0, 4.5, 5.0, 3.5 & MW \\
$R_{linea}$ (L1..L4) & 0.04, 0.06, 0.05, 0.07 & pu \\
$X_{linea}$ (L1..L4) & 0.08, 0.10, 0.09, 0.11 & pu \\
$T_{max,dev}$ & 65 & \si{\degree}C} \\
\bottomrule
\end{tabular}
\end{table}

\subsection{Catalogo de fallas}

El catalogo incluye ocho tipos de falla representativos del SEN
chileno segun statistic de SEC (2024) y experiencia de
distribuidoras~\citep{SEC-Estadisticas-2024}. Las fallas agudas
(F1-F3) tienen distribucion de Poisson; las lentas (F4, F6, F7)
distribucion Weibull; las climaticas (F8) Poisson anual.

\begin{table}[ht]
\centering
\small
\caption{Catalogo de fallas sinteticas.}
\label{tab:catalogofallas}
\begin{tabular}{lccr}
\toprule
\textbf{Codigo} & \textbf{Tipo} & \textbf{Severidad} & $\lambda$ / dia \\
\midrule
F1 & Cortocircuito monofasico L1   & 0.85 & 0.50 \\
F2 & Sobrecarga sostenida L2       & 0.45 & 0.30 \\
F3 & Sobretension AT               & 0.90 & 0.10 \\
F4 & Degradacion de aislamiento    & 0.30 & 0.40 \\
F5 & Perdida de comunicacion PMU   & 0.50 & 0.35 \\
F6 & Desbalance de cargas          & 0.35 & 0.60 \\
F7 & Falla incipiente de buje      & 0.25 & 0.20 \\
F8 & Evento climatico extremo      & 0.70 & $5/365$ \\
\bottomrule
\end{tabular}
\end{table}

\subsection{Perfil de carga y clima}

El perfil de carga horario modelado corresponde a una distribuidora
chilena tipica: valle en la madrugada (02:00-05:00), pico en la noche
(19:00-22:00) por consumo residencial, y un pico de mediodia mas
suave por consumo comercial. La formula utilizada es:

\[
f(t) = 0.55 + 0.30\, e^{-\frac{(h - 20)^2}{8}} + 0.15\, e^{-\frac{(h - 12)^2}{4.5}}
\]

mas ruido gaussiano $\mathcal{N}(0, 0.04)$ para la variabilidad
intra-hora. El perfil es valido para una SE urbana mixta
(residencial + comercial ligero); las SE rurales tienen un pico
agricola al amanecer y al atardecer que no se modela aqui pero
puede agregarse parametrizando la curva.

El clima extremo (F8) se modela como un frente que produce una
caida simultanea de tension AT ($V_{AT} - 0.06 \cdot sev$), una
desviacion de frecuencia ($-0.12 \cdot sev$~Hz) y un aumento
sostenido de corriente en todos los alimentadores (las cargas
intentan compensar la caida de tension).

\subsection{Cadencia y dimensionalidad}

\begin{itemize}[leftmargin=*,itemsep=2pt]
  \item \textbf{SCADA} a 1~Hz: 5 tensiones, 4 corrientes, 2 potencias,
        2 temperaturas, THD V/I, angulos.
  \item \textbf{PMU} a 20~Hz: tension AT, corriente AT, frecuencia,
        ROCOF.
  \item En este trabajo se usan solo las observaciones a 1~Hz para
        simplicidad; el pipeline escala trivialmente a 20~Hz agregando
        capas de downsample.
\end{itemize}

El vector de observaciones por instante es de 22 features. Las
ventanas del autoencoder LSTM son de 30 muestras (30~s), lo que da
un tensor de entrada de $30 \times 22$.

\subsection{Conjuntos de entrenamiento, validacion y test}

Se simulan 7 dias ($\approx 604{,}800$ muestras) y se particionan:

\begin{itemize}[leftmargin=*,itemsep=2pt]
  \item Train: primeras 6 muestras (70\%); de este subconjunto,
        \emph{solo} las ventanas sin falla se usan para entrenar el
        autoencoder.
  \item Validacion: 15\% siguientes, para early stopping y calibracion
        conformal.
  \item Test: 15\% finales, para evaluar metricas.
\end{itemize}

La particion se hace al nivel de ventana y es estratificada por
etiqueta de falla para mantener la proporcion de positivos.

\subsection{Sanidad de los datos sinteticos}

Para verificar que el gemelo produce dinamica realista, se compararon
las estadisticas de los residuos en regimen normal contra una muestra
de medidas de frecuencia del CEN (Coordinador Electrico Nacional,
datos publicos). El match es razonablemente bueno:

\begin{table}[ht]
\centering
\small
\caption{Comparacion de dinamica de fondo: gemelo vs CEN.}
\label{tab:cenmatch}
\begin{tabular}{lcc}
\toprule
\textbf{Estadistica} & \textbf{Gemelo} & \textbf{CEN (real)} \\
\midrule
$\mu(f)$   & 50.00 Hz  & 50.02 Hz \\
$\sigma(f)$ & 0.045 Hz & 0.062 Hz \\
$\mu(ROCOF)$ & 0.00 Hz/s & $-0.001$ Hz/s \\
$\sigma(ROCOF)$ & 0.020 Hz/s & 0.025 Hz/s \\
\bottomrule
\end{tabular}
\end{table}

Las pequenas diferencias son consistentes con la hipotesis de que el
gemelo es un modelo de bajo orden. Para uso en produccion, el gemelo
se recalibra con los primeros 30 dias de operacion de la SE real
usando un estimador de maxima verosimilitud sobre los parametros
sensibles ($R$, $X$, $P_{cu}$).